\section{Polymer geometry}\label{app:polymer-geometry}

We use the geometric arguments of \cite[Appendix~A]{QED3Manuscript}.
Let $L$ be odd, $M$ a power of $L$, and $\ell$ a multiple of $LM$
with $\ell\ge13LM$. The $M$-cubes have centres in
$M\Z^3/\ell\Z^3$. They partition the lattice sites; their closed
continuum versions may meet along faces and edges. A polymer is a
nonempty union of $M$-cubes connected through faces. Denote its number
of cubes by $|X|_M$. A net in the continuum torus is a connected
finite union of images of specified closed straight segments
in $\R^3$, a point being allowed as a segment.
Its length is the sum of the Euclidean lengths of those segments.
Set
\[
 Md_M(X)=\inf\{\operatorname{length}(N):N
             \text{ is a net meeting every closed cube of }X\},
\]
where the infimum also extends over the finite segment representations.
The same definition is used for finite face-connected unions of
$M$-cubes in $\R^3$, with nets in $\R^3$. Distances are sup distances.


\subsection{Size and diameter}
The following bounds hold on the torus and on the covering space:
\begin{align}
 d_M(X)&\le |X|_M-1,&
 |X|_M&\le54(d_M(X)+1),\label{eq:polymer-size-tree}\\
 d(p,q)&\le M(d_M(X)+2)
       &&(p,q\text{ in the closed cubes of }X).
                                  \label{eq:polymer-diameter}
\end{align}

To prove the first inequality of \eqref{eq:polymer-size-tree}, order
the cubes along a spanning tree of their adjacency graph. A net
meeting one cube can be extended to an adjacent cube by a segment of
length at most $M$ perpendicular to their common face. Starting with
one point and repeating this construction gives a net of length at
most $M(|X|_M-1)$.

For the reverse inequality, subdivide the segments of any net $N$ at
their intersections and discard overlapping pieces. Its total length
does not increase. Doubling the edges of the resulting connected
graph gives a closed traversal of length at most
$2\operatorname{length}(N)$. Cut it into at most
$2\operatorname{length}(N)/M+1$ arcs of length at most $M$.
Each arc lifts to a curve in $\R^3$ whose range in each coordinate
has length at most $M$. Such an interval meets at most three closed
cube intervals, so an arc meets at most $27$ closed $M$-cubes.
Every cube of $X$ meets one of these arcs. Taking the infimum over
$N$ gives
\[
 |X|_M\le27(2d_M(X)+1)\le c_X(d_M(X)+1).
\]
For \eqref{eq:polymer-diameter}, take a net of length less than
$Md_M(X)+\delta$ and points $p',q'$ of the net in the respective
cubes of $p,q$. The path in the net from $p'$ to $q'$ has length
at most that of the net. Hence
$d(p,q)\le M+(Md_M(X)+\delta)+M$; let $\delta\downarrow0$.
The same argument holds in $\R^3$.

\subsection{Nonwrapping lifts}\label{subsec:polymer-lifts}

Let $\operatorname{pr}:\R^3\to\R^3/\ell\Z^3$ be the projection.
\begin{proposition}\label{prop:nonwrapping-polymer-lift}
If $X$ is a polymer and
\begin{equation}\label{eq:nonwrapping-polymer-condition}
 2M(d_M(X)+2)<\ell,
\end{equation}
then $X=\operatorname{pr}(\widetilde X)$ for a finite face-connected
union of $M$-cubes in $\R^3$. The projection is injective on
$\widetilde X$, preserves adjacency of its cubes in both directions,
and $d_M(\widetilde X)=d_M(X)$. The lift is unique up to a translation
in $\ell\Z^3$.
\end{proposition}

\begin{proof}
Put $D=M(d_M(X)+2)$, fix $p_0$ in a closed cube of $X$, and fix its
lift $\widetilde p_0$. The projection is injective on the closed sup
ball $\widetilde B=\{|p-\widetilde p_0|_\infty\le D\}$, since two
points there differ in every coordinate by at most $2D<\ell$.
It maps this ball onto the torus ball $B=\{p:d(p,p_0)\le D\}$.
Its inverse $\varphi:B\to\widetilde B$ is continuous, and is locally
isometric. The diameter estimate places every closed cube of $X$ in
$B$. A segment or a closed $M$-cube lying in $B$ lifts under
$\varphi$ to a segment of the same length or a closed $M$-cube:
the translates $\widetilde B+\ell n$ are pairwise disjoint, and a
connected lift lies in one of them. The same inverse maps shared
faces to shared faces. Thus $\widetilde X=\varphi(X)$ has the
asserted injectivity, adjacency, and connectedness.

Any other face-connected injective lift can be translated so that it
contains the chosen lift of the cube of $p_0$. There are no two
different neighbours of one lifted cube with the same projection:
their centres differ by at most $2M<\ell$. Propagate this observation
along a chain of adjacent cubes. Every cube of the other lift equals
the corresponding cube of $\widetilde X$, proving uniqueness.

Projecting a net for $\widetilde X$ gives a net for $X$ with no greater
length, so $d_M(X)\le d_M(\widetilde X)$. Conversely, let $N$ be a
net for $X$ of length at most $Md_M(X)+M$, and choose its point $q_0$
in the cube of $p_0$. Every point of $N$ is joined to $q_0$ by a
curve of length at most that of $N$, and $d(q_0,p_0)\le M$.
Therefore $N\subset B$. Its lift $\varphi(N)$ has the same length
and meets every cube of $\widetilde X$. Taking an infimizing sequence
of these nets proves the reverse inequality.
\end{proof}

For the determinant estimate, suppose in addition that
$M(d_M(X)+4)<\ell$ and $M\ge2$. The closed cubes of $\widetilde X$
have coordinate ranges of length at most $M(d_M(X)+2)$ by
\eqref{eq:polymer-diameter}. Enlarge these intervals by one mesh
$\eta_k\le1$ at each end and intersect with $\eta_k\Z^3$. The
resulting box $W$ contains every bond meeting the lifted polymer and
has each side of length at most $M(d_M(X)+3)$. This length, increased
by one mesh, is less than $\ell$. Thus the projection is injective
on $W$, and every torus bond meeting $X$ has exactly one lift in $W$.

